\section{Existence of the blowup along a smooth centre}
\label{sec:blowup-existence}

This section constructs the blowup of a complex manifold \(X\) of dimension \(n\) (Hausdorff, second countable) along a nowhere dense
closed submanifold \(C\) in the sense of Definition~\ref{ana:blowups}, and so proves Lemma~\ref{ana:blowup-existence}. The construction
is Demailly's, \emph{Complex analytic and differential geometry}, VII.12: as a set the blowup is \((X\setminus C)\sqcup P(NC)\), and its
manifold structure is given by an explicit atlas whose coordinate changes are holomorphic because of a division lemma of Hadamard type.
The route differs from the source as follows.
\begin{enumerate}
\item The components of \(C\) may have different dimensions (Definition~\ref{ana:blowups} only asks that \(C\) be regular of some dimension at
each of its points); nothing below uses a fixed codimension.
\item The normal coordinates are the \emph{last} \(n-k\) coordinates: a flattening chart maps \(C\) onto \(\mathbf C^k\times0\) (Demailly:
\(C=\{z_1=\dots=z_s=0\}\)), and flattening charts are arbitrary biholomorphisms of open sets, not products of balls.
\item No normal bundle is constructed. The tangent space \(T_xX\) is identified with \(\mathbf C^n\) by the preferred chart at \(x\), the tangent
space \(T_xC\) is defined as the common kernel of the differentials of the germs of the ideal of \(C\) at \(x\), and a point of \(P(N_xC)\) is
recorded as the subspace \(T_xC+\mathbf C\xi\) of \(T_xX\) (\(\xi\notin T_xC\)).
\item The universal property of Definition~\ref{ana:blowups} is not part of Demailly's definition (12.3); it is proved here
(Lemmas~\ref{bex:lift-unique} and~\ref{bex:lift-local}).
\end{enumerate}

\paragraph{Notation.} For \(x\in X\) let \(e_x\) be the preferred chart at \(x\). For a function \(f\) holomorphic near \(x\) let
\(d_xf:=D(f\circ e_x^{-1})(e_x(x))\colon\mathbf C^n\to\mathbf C\), and let \(\mathcal I_{C,x}\) be the set of functions holomorphic near \(x\) that
vanish on \(C\) near \(x\). The \emph{tangent space} of \(C\) at \(x\) is
\[
 T_xC:=\bigcap_{f\in\mathcal I_{C,x}}\ker d_xf\subset\mathbf C^n .
\]
A \emph{flattening chart} \(c=(y_0,k,\Phi,j)\) consists of a point \(y_0\in X\), a biholomorphism \(\Phi\) between open subsets of \(\mathbf C^n\)
with \(\Phi.\mathrm{source}\subset e_{y_0}.\mathrm{target}\) such that, for \(x\) in the \emph{flat domain}
\(D_\Phi=\{x\in e_{y_0}.\mathrm{source}:e_{y_0}(x)\in\Phi.\mathrm{source}\}\), \(x\in C\) if and only if \(\Phi(e_{y_0}(x))\in\mathbf C^k\times0\),
and a \emph{normal index} \(j\) with \(k\le j<n\). Its \emph{flat coordinates} are \(\tau_i=(\Phi\circ e_{y_0})_i\) on \(D_\Phi\);
\(\tau_i\in\mathcal I_{C,x}\) for \(x\in D_\Phi\) and \(i\ge k\). Since \(C\) is a closed submanifold, every point of \(C\) lies in the flat domain of
a flattening (Definition~\ref{ana:blowups}). We write \(\sigma_j\colon\mathbf C^n\to\mathbf C^n\) for the monomial map
\(\sigma_j(w)_i=w_iw_j\) (\(k\le i\ne j\)), \(\sigma_j(w)_i=w_i\) otherwise.

\begin{lemma}[Hadamard's lemma in flattening coordinates]
\label{bex:hadamard}
\uses{ana:blowups}
\lean{TranscendentalNakamaye.Pending.CoordinatePlaneDivision.exists_eq_coord_mul,TranscendentalNakamaye.Pending.CoordinatePlaneDivision.exists_eq_sum_coord_mul,TranscendentalNakamaye.Pending.BlowupFlatCoordinates.flatCoord,TranscendentalNakamaye.Pending.BlowupFlatCoordinates.flatDomain,TranscendentalNakamaye.Pending.BlowupFlatCoordinates.holomorphicAt_flatCoord,TranscendentalNakamaye.Pending.BlowupFlatCoordinates.flatCoord_mem_vanishingIdeal,TranscendentalNakamaye.Pending.BlowupFlatCoordinates.exists_eq_sum_flatCoord}
\leanok
\begin{enumerate}
\item Let \(F\) be holomorphic near a point \(q\in\mathbf C^k\times0\subset\mathbf C^n\) and vanish on \(\mathbf C^k\times0\) near \(q\). Then
\(F(u)=\sum_{i\ge k}u_iH_i(u)\) near \(q\) with \(H_i\) holomorphic near \(q\).
\item For a flattening chart and \(x\in D_\Phi\), every \(f\in\mathcal I_{C,x}\) is \(\sum_{i\ge k}b_i\tau_i\) near \(x\) with \(b_i\) holomorphic near \(x\).
\end{enumerate}
\end{lemma}
\begin{proof}
\leanok
(1) Induction on the number of normal coordinates: \(F(u)-F(u_1,\dots,u_{n-1},0)=u_n\int_0^1\partial_nF(u_1,\dots,u_{n-1},tu_n)\,dt\).
(2) Apply (1) to \(F=f\circ(\Phi\circ e_{y_0})^{-1}\) at \(\Phi(e_{y_0}(x))\) if \(x\in C\); if \(x\notin C\) some \(\tau_i(x)\ne0\) and
\(f=(f/\tau_i)\,\tau_i\) near \(x\).
\end{proof}

\begin{lemma}[The tangent space in flattening coordinates]
\label{bex:tangent-flat}
\uses{bex:hadamard}
\lean{TranscendentalNakamaye.Pending.BlowupSpaceTangent.differentiableAt_comp_chart_symm,TranscendentalNakamaye.Pending.BlowupSpaceTangent.differentiableAt_chart_comp_chart_symm,TranscendentalNakamaye.Pending.BlowupSpaceTangent.germDiff,TranscendentalNakamaye.Pending.BlowupSpaceTangent.chartAt_symm_self,TranscendentalNakamaye.Pending.BlowupSpaceTangent.differentiableAt_germ,TranscendentalNakamaye.Pending.BlowupSpaceTangent.tendsto_chartAt_symm,TranscendentalNakamaye.Pending.BlowupSpaceTangent.germDiff_congr,TranscendentalNakamaye.Pending.BlowupSpaceTangent.germDiff_zero,TranscendentalNakamaye.Pending.BlowupSpaceTangent.germDiff_mul,TranscendentalNakamaye.Pending.BlowupSpaceTangent.germDiff_sum,TranscendentalNakamaye.Pending.BlowupSpaceTangent.tangentSpace,TranscendentalNakamaye.Pending.BlowupSpaceTangent.mem_tangentSpace,TranscendentalNakamaye.Pending.BlowupSpaceTangent.germDiff_eq_zero_of_mem_tangentSpace,TranscendentalNakamaye.Pending.BlowupSpaceTangent.flatMap,TranscendentalNakamaye.Pending.BlowupSpaceTangent.flatInvMap,TranscendentalNakamaye.Pending.BlowupSpaceTangent.flatVec,TranscendentalNakamaye.Pending.BlowupSpaceTangent.flatInvMap_flat,TranscendentalNakamaye.Pending.BlowupSpaceTangent.flatInverse_flatMap_eventually,TranscendentalNakamaye.Pending.BlowupSpaceTangent.differentiableAt_flatMap,TranscendentalNakamaye.Pending.BlowupSpaceTangent.differentiableAt_flatInvMap,TranscendentalNakamaye.Pending.BlowupSpaceTangent.flatMap_flatInvMap_eventually,TranscendentalNakamaye.Pending.BlowupSpaceTangent.fderiv_flatMap_flatVec,TranscendentalNakamaye.Pending.BlowupSpaceTangent.germDiff_flatCoord,TranscendentalNakamaye.Pending.BlowupSpaceTangent.germDiff_flatCoord_flatVec,TranscendentalNakamaye.Pending.BlowupSpaceTangent.germDiff_flatVec,TranscendentalNakamaye.Pending.BlowupSpaceTangent.flatCoord_eq_zero,TranscendentalNakamaye.Pending.BlowupSpaceTangent.mem_tangentSpace_iff,TranscendentalNakamaye.Pending.BlowupSpaceTangent.flatVec_not_mem_tangentSpace,TranscendentalNakamaye.Pending.BlowupSpaceTangent.exists_sub_smul_mem_of_mem_sup_span,TranscendentalNakamaye.Pending.BlowupSpaceTangent.sup_span_eq_of_sub_smul_mem,TranscendentalNakamaye.Pending.BlowupSpaceTangent.exists_sub_smul_mem_of_sup_span_eq,TranscendentalNakamaye.Pending.BlowupSpaceTangent.germDiff_apply_of_sub_smul_mem}
\leanok
Let \(c\) be a flattening chart and \(x\in C\cap D_\Phi\).
\begin{enumerate}
\item \(d_xf\) vanishes on \(T_xC\) for every \(f\in\mathcal I_{C,x}\).
\item \(T_xC=\{\xi\in\mathbf C^n:d_x\tau_i(\xi)=0\text{ for all }i\ge k\}\).
\item The linear map \(\xi\mapsto(d_x\tau_i(\xi))_{i}\) is an isomorphism \(\mathbf C^n\to\mathbf C^n\).
\end{enumerate}
\end{lemma}
\begin{proof}
\leanok
(1) is the definition of \(T_xC\). In (2), ``\(\subset\)'' holds because \(\tau_i\in\mathcal I_{C,x}\) for \(i\ge k\). Conversely, for
\(f\in\mathcal I_{C,x}\) write \(f=\sum_{i\ge k}b_i\tau_i\) near \(x\) (Lemma~\ref{bex:hadamard}); since \(\tau_i(x)=0\) for \(i\ge k\) (as \(x\in C\)), the
product rule gives \(d_xf=\sum_{i\ge k}b_i(x)\,d_x\tau_i\), which vanishes wherever all \(d_x\tau_i\), \(i\ge k\), do. (3) The map is the
differential at \(e_x(x)\) of \(\Phi\circ e_{y_0}\circ e_x^{-1}\), which has the holomorphic local inverse \(e_x\circ e_{y_0}^{-1}\circ\Phi^{-1}\).
\end{proof}

\begin{definition}[The blowup as a set]
\label{bex:blowup-space}
\uses{bex:tangent-flat}
\source{demailly-cadg:VII.12.1}
\lean{TranscendentalNakamaye.Pending.BlowupConstruction.IsBlowupPoint,TranscendentalNakamaye.Pending.BlowupConstruction.BlowupSpace,TranscendentalNakamaye.Pending.BlowupConstruction.BlowupSpace.base,TranscendentalNakamaye.Pending.BlowupConstruction.BlowupSpace.dir,TranscendentalNakamaye.Pending.BlowupConstruction.BlowupSpace.ext,TranscendentalNakamaye.Pending.BlowupConstruction.BlowupSpace.isBlowupPoint,TranscendentalNakamaye.Pending.BlowupConstruction.BlowupSpace.dir_eq_bot,TranscendentalNakamaye.Pending.BlowupConstruction.BlowupSpace.exists_dir,TranscendentalNakamaye.Pending.BlowupConstruction.BlowupSpace.dirVec,TranscendentalNakamaye.Pending.BlowupConstruction.BlowupSpace.dirVec_not_mem,TranscendentalNakamaye.Pending.BlowupConstruction.BlowupSpace.dir_eq_sup_dirVec,TranscendentalNakamaye.Pending.BlowupConstruction.blowupProj,TranscendentalNakamaye.Pending.BlowupConstruction.offPoint,TranscendentalNakamaye.Pending.BlowupConstruction.excPoint,TranscendentalNakamaye.Pending.BlowupConstruction.eq_offPoint,TranscendentalNakamaye.Pending.BlowupConstruction.eq_of_base_eq_of_not_mem,TranscendentalNakamaye.Pending.BlowupConstruction.extVal,TranscendentalNakamaye.Pending.BlowupConstruction.extVal_of_not_mem,TranscendentalNakamaye.Pending.BlowupConstruction.extVal_of_mem,TranscendentalNakamaye.Pending.BlowupConstruction.exists_extVal_eq_mul,TranscendentalNakamaye.Pending.BlowupConstruction.exists_extVal_excPoint_eq_mul}
\leanok
Let \(\tilde X\) be the set of pairs \((x,L)\), \(x\in X\) and \(L\subset\mathbf C^n\) a linear subspace, such that either \(x\notin C\) and
\(L=0\), or \(x\in C\) and \(L=T_xC+\mathbf C\xi\) for some \(\xi\notin T_xC\); let \(\pi(x,L)=x\). For \(p=(x,L)\) with \(x\in C\) fix
\(\xi_p\notin T_xC\) with \(L=T_xC+\mathbf C\xi_p\). For \(f\) holomorphic near \(x\) put \(f^\flat(p)=f(x)\) if \(x\notin C\) and
\(f^\flat(p)=d_xf(\xi_p)\) if \(x\in C\).
\end{definition}
If \(T_xC+\mathbf C\xi=T_xC+\mathbf C\xi'\) with \(\xi,\xi'\notin T_xC\), then \(\xi'=a\xi+t\) with \(a\ne0\) and \(t\in T_xC\); so for
\(f,g\in\mathcal I_{C,x}\) the quotient \(f^\flat(p)/g^\flat(p)\) does not depend on the choice of \(\xi_p\) (Lemma~\ref{bex:tangent-flat}(1)). It is
Demailly's extension \((f/g)^\sim\) of \(f/g\). Conversely, by Lemma~\ref{bex:tangent-flat}(2), if \(x\in C\cap D_\Phi\) and
\((d_x\tau_i(\xi'))_{i\ge k}=a\,(d_x\tau_i(\xi))_{i\ge k}\) with \(a\ne0\), then \(\xi'-a\xi\in T_xC\) and \(T_xC+\mathbf C\xi=T_xC+\mathbf C\xi'\).

\begin{lemma}[The monomial charts]
\label{bex:monomial-chart}
\uses{bex:blowup-space,bex:tangent-flat}
\source{demailly-cadg:VII.12.1,demailly-cadg:VII.12.2}
\lean{TranscendentalNakamaye.Pending.BlowupConstruction.FlatChart,TranscendentalNakamaye.Pending.BlowupConstruction.FlatChart.τ,TranscendentalNakamaye.Pending.BlowupConstruction.FlatChart.dom,TranscendentalNakamaye.Pending.BlowupConstruction.FlatChart.τ_eq_zero,TranscendentalNakamaye.Pending.BlowupConstruction.FlatChart.τ_mem_vanishingIdeal,TranscendentalNakamaye.Pending.BlowupConstruction.FlatChart.pt,TranscendentalNakamaye.Pending.BlowupConstruction.FlatChart.target,TranscendentalNakamaye.Pending.BlowupConstruction.FlatChart.isOpen_target,TranscendentalNakamaye.Pending.BlowupConstruction.FlatChart.pt_mem_dom,TranscendentalNakamaye.Pending.BlowupConstruction.FlatChart.flat_pt,TranscendentalNakamaye.Pending.BlowupConstruction.FlatChart.τ_pt,TranscendentalNakamaye.Pending.BlowupConstruction.FlatChart.pt_mem_iff,TranscendentalNakamaye.Pending.BlowupConstruction.FlatChart.normalVec,TranscendentalNakamaye.Pending.BlowupConstruction.FlatChart.normalVec_j,TranscendentalNakamaye.Pending.BlowupConstruction.FlatChart.junk,TranscendentalNakamaye.Pending.BlowupConstruction.FlatChart.invFun,TranscendentalNakamaye.Pending.BlowupConstruction.FlatChart.toFun,TranscendentalNakamaye.Pending.BlowupConstruction.FlatChart.source,TranscendentalNakamaye.Pending.BlowupConstruction.FlatChart.invFun_of_ne,TranscendentalNakamaye.Pending.BlowupConstruction.FlatChart.flatVec_normalVec_not_mem,TranscendentalNakamaye.Pending.BlowupConstruction.FlatChart.invFun_of_eq,TranscendentalNakamaye.Pending.BlowupConstruction.FlatChart.base_invFun,TranscendentalNakamaye.Pending.BlowupConstruction.FlatChart.monomial_toFun,TranscendentalNakamaye.Pending.BlowupConstruction.FlatChart.eq_excPoint,TranscendentalNakamaye.Pending.BlowupConstruction.FlatChart.toFun_apply_j,TranscendentalNakamaye.Pending.BlowupConstruction.FlatChart.toFun_apply_normal,TranscendentalNakamaye.Pending.BlowupConstruction.FlatChart.toFun_apply_of_not,TranscendentalNakamaye.Pending.BlowupConstruction.FlatChart.map_source,TranscendentalNakamaye.Pending.BlowupConstruction.FlatChart.pt_toFun,TranscendentalNakamaye.Pending.BlowupConstruction.FlatChart.toFun_j_eq_zero_iff,TranscendentalNakamaye.Pending.BlowupConstruction.FlatChart.monomial_apply_of_not,TranscendentalNakamaye.Pending.BlowupConstruction.FlatChart.map_target,TranscendentalNakamaye.Pending.BlowupConstruction.FlatChart.left_inv,TranscendentalNakamaye.Pending.BlowupConstruction.FlatChart.right_inv,TranscendentalNakamaye.Pending.BlowupConstruction.FlatChart.toPartialEquiv,TranscendentalNakamaye.Pending.BlowupConstruction.FlatChart.blowupProj_invFun,TranscendentalNakamaye.Pending.BlowupConstruction.OffChart,TranscendentalNakamaye.Pending.BlowupConstruction.OffChart.e,TranscendentalNakamaye.Pending.BlowupConstruction.OffChart.toPartialEquiv,TranscendentalNakamaye.Pending.BlowupConstruction.OffChart.source_eq,TranscendentalNakamaye.Pending.BlowupConstruction.OffChart.target_eq,TranscendentalNakamaye.Pending.BlowupConstruction.OffChart.apply_eq,TranscendentalNakamaye.Pending.BlowupConstruction.OffChart.symm_apply_of_not_mem}
\leanok
Let \(c=(y_0,k,\Phi,j)\) be a flattening chart, \(U_c=\{p\in\tilde X:\pi(p)\in D_\Phi,\ \tau_j^\flat(p)\ne0\}\), and
\(\chi_c(p)=w\) with \(w_i=\tau_i(\pi(p))\) for \(i<k\) or \(i=j\), and \(w_i=\tau_i^\flat(p)/\tau_j^\flat(p)\) for \(k\le i\ne j\). Then
\(\Phi(e_{y_0}(\pi(p)))=\sigma_j(\chi_c(p))\), and \(\chi_c\) is a bijection of \(U_c\) onto \(\sigma_j^{-1}(\Phi.\mathrm{target})\). Its inverse
sends \(w\) to \((x,0)\) if \(w_j\ne0\) and to \((x,T_xC+\mathbf C\xi_w)\) if \(w_j=0\), where \(x=(\Phi\circ e_{y_0})^{-1}(\sigma_j(w))\) and
\(\xi_w\) is the vector with \(d_x\tau_i(\xi_w)=v_i\) for all \(i\) (Lemma~\ref{bex:tangent-flat}(3)), \(v_i=0\) for \(i<k\), \(v_j=1\), \(v_i=w_i\) for
\(k\le i\ne j\). Similarly, for a chart \(e\) of \(X\), \(\chi_e(p)=e(\pi(p))\) is a bijection of
\(U_e=\pi^{-1}(e.\mathrm{source}\setminus C)\) onto \(e(e.\mathrm{source}\setminus C)\) with inverse \(z\mapsto(e^{-1}(z),0)\).
\end{lemma}
\begin{proof}
\leanok
Let \(p\in U_c\), \(x=\pi(p)\), \(w=\chi_c(p)\). If \(x\notin C\) then \(\tau^\flat=\tau\) at \(p\), \(w_j=\tau_j(x)\ne0\) and
\(w_iw_j=\tau_i(x)\) for \(k\le i\ne j\). If \(x\in C\) then \(\tau_i(x)=0\) for \(i\ge k\), \(w_j=0\) and \(\sigma_j(w)_i=w_iw_j=0\) for
\(k\le i\ne j\). In both cases \(\sigma_j(w)=\Phi(e_{y_0}(x))\in\Phi.\mathrm{target}\), so \(x\) is recovered from \(w\). If \(w_j\ne0\) then
\(\sigma_j(w)\notin\mathbf C^k\times0\), so \(x\notin C\) and \(p=(x,0)\). If \(w_j=0\) then \(x\in C\), and
\((d_x\tau_i(\xi_p))_{i\ge k}=\tau_j^\flat(p)\,(v_i)_{i\ge k}=\tau_j^\flat(p)\,(d_x\tau_i(\xi_w))_{i\ge k}\), so \(p=(x,T_xC+\mathbf C\xi_w)\) by the
remark after Definition~\ref{bex:blowup-space}. Conversely, for \(w\in\sigma_j^{-1}(\Phi.\mathrm{target})\) the pair described is a point of
\(\tilde X\) (for \(w_j=0\): \(\sigma_j(w)\in\mathbf C^k\times0\), so \(x\in C\), and \(\xi_w\notin T_xC\) because \(v_j=1\)); it lies in \(U_c\), as
\(\tau_j^\flat=\tau_j(x)=w_j\ne0\), resp.\ \(\tau_j^\flat=a\,v_j=a\ne0\) when \(\xi_p=a\xi_w+t\); and \(\chi_c\) maps it to \(w\), the ratios being
\(a\,v_i/a\,v_j=w_i\). The statement about \(\chi_e\) is clear.
\end{proof}

\begin{lemma}[Extended quotients are holomorphic in the monomial charts]
\label{bex:ratio-holomorphic}
\uses{bex:monomial-chart,bex:hadamard}
\source{demailly-cadg:VII.12.2}
\lean{TranscendentalNakamaye.Pending.BlowupConstruction.FlatChart.ratio,TranscendentalNakamaye.Pending.BlowupConstruction.FlatChart.flatRep,TranscendentalNakamaye.Pending.BlowupConstruction.FlatChart.flatRep_monomial,TranscendentalNakamaye.Pending.BlowupConstruction.FlatChart.monomial_apply_j,TranscendentalNakamaye.Pending.BlowupConstruction.FlatChart.ratio_of_ne,TranscendentalNakamaye.Pending.BlowupConstruction.FlatChart.ratio_of_eq,TranscendentalNakamaye.Pending.BlowupConstruction.FlatChart.differentiableAt_flatRep,TranscendentalNakamaye.Pending.BlowupConstruction.FlatChart.differentiableAt_comp_pt,TranscendentalNakamaye.Pending.BlowupConstruction.FlatChart.continuousAt_pt,TranscendentalNakamaye.Pending.BlowupConstruction.FlatChart.hadamardExpr,TranscendentalNakamaye.Pending.BlowupConstruction.FlatChart.normalVec_apply,TranscendentalNakamaye.Pending.BlowupConstruction.FlatChart.monomial_mul_eq,TranscendentalNakamaye.Pending.BlowupConstruction.FlatChart.fderiv_sum_coord_mul,TranscendentalNakamaye.Pending.BlowupConstruction.FlatChart.differentiableAt_ratio}
\leanok
Let \(c=(y_0,k,\Phi,j)\) be a flattening chart, \(V\subset X\) open, and \(f\) holomorphic on \(V\) and zero on \(C\cap V\). Then
\(G_f(w)=f^\flat(\chi_c^{-1}(w))/\tau_j^\flat(\chi_c^{-1}(w))\) is holomorphic on the open set
\(\{w\in\sigma_j^{-1}(\Phi.\mathrm{target}):(\Phi\circ e_{y_0})^{-1}(\sigma_j(w))\in V\}\).
\end{lemma}
\begin{proof}
\leanok
Let \(w_0\) be in this set and \(u_0=\sigma_j(w_0)\). If \((w_0)_j\ne0\), then near \(w_0\) the points \(\chi_c^{-1}(w)\) lie off \(C\) and
\(G_f(w)=f((\Phi\circ e_{y_0})^{-1}(\sigma_j(w)))/w_j\). Let \((w_0)_j=0\), so \(u_0\in\mathbf C^k\times0\). The function
\(F=f\circ(\Phi\circ e_{y_0})^{-1}\) is holomorphic near \(u_0\) and vanishes on \(\mathbf C^k\times0\) near \(u_0\); by
Lemma~\ref{bex:hadamard}, \(F(u)=\sum_{i\ge k}u_iH_i(u)\) near \(u_0\). For \(w\) near \(w_0\) with \(w_j\ne0\),
\[
 G_f(w)=\frac{F(\sigma_j(w))}{w_j}=H_j(\sigma_j(w))+\sum_{k\le i\ne j}w_i\,H_i(\sigma_j(w)).
\]
For \(w\) near \(w_0\) with \(w_j=0\), let \(x=(\Phi\circ e_{y_0})^{-1}(\sigma_j(w))\in C\) and \(\xi_w\) be as in Lemma~\ref{bex:monomial-chart}; since
\(\xi_p=a\xi_w+t\) (\(a\ne0\), \(t\in T_xC\)) and \(d_x\tau_j(\xi_w)=1\), \(G_f(w)=d_xf(\xi_w)\). Near \(e_x(x)\),
\(f\circ e_x^{-1}=F\circ(\Phi\circ e_{y_0}\circ e_x^{-1})\), and \(D(\Phi\circ e_{y_0}\circ e_x^{-1})(\xi_w)=v\); so, with \(u=\sigma_j(w)\),
\[
 G_f(w)=DF(u)(v)=\sum_{i\ge k}v_iH_i(u)+\sum_{i\ge k}u_i\,DH_i(u)(v)=H_j(\sigma_j(w))+\sum_{k\le i\ne j}w_i\,H_i(\sigma_j(w)),
\]
because \(u_i=0\) for \(i\ge k\). So \(G_f\) is given near \(w_0\) by a single holomorphic expression.
\end{proof}

\begin{theorem}[The blowup is a complex manifold]
\label{bex:atlas}
\uses{bex:monomial-chart,bex:ratio-holomorphic,bex:tangent-flat}
\source{demailly-cadg:VII.12.2,demailly-cadg:VII.12.3}
\lean{TranscendentalNakamaye.Pending.BlowupConstruction.blowupAtlas,TranscendentalNakamaye.Pending.BlowupConstruction.FlatChart.toPartialEquiv_source,TranscendentalNakamaye.Pending.BlowupConstruction.FlatChart.toPartialEquiv_target,TranscendentalNakamaye.Pending.BlowupConstruction.FlatChart.toPartialEquiv_apply,TranscendentalNakamaye.Pending.BlowupConstruction.FlatChart.toPartialEquiv_symm_apply,TranscendentalNakamaye.Pending.BlowupConstruction.FlatChart.continuousOn_pt,TranscendentalNakamaye.Pending.BlowupConstruction.FlatChart.isOpen_target_inter_pt,TranscendentalNakamaye.Pending.BlowupConstruction.FlatChart.isOpen_dom,TranscendentalNakamaye.Pending.BlowupConstruction.FlatChart.holomorphicAt_τ,TranscendentalNakamaye.Pending.BlowupConstruction.FlatChart.mem_source_iff,TranscendentalNakamaye.Pending.BlowupConstruction.FlatChart.extVal_invFun_ne,TranscendentalNakamaye.Pending.BlowupConstruction.FlatChart.differentiableAt_ratio_τ,TranscendentalNakamaye.Pending.BlowupConstruction.FlatChart.continuousOn_ratio_τ,TranscendentalNakamaye.Pending.BlowupConstruction.OffChart.toPartialEquiv_symm_apply_of_not_mem,TranscendentalNakamaye.Pending.BlowupConstruction.transSource,TranscendentalNakamaye.Pending.BlowupConstruction.transSource_flat_flat,TranscendentalNakamaye.Pending.BlowupConstruction.isOpen_transSource_flat_flat,TranscendentalNakamaye.Pending.BlowupConstruction.differentiableOn_flat_flat,TranscendentalNakamaye.Pending.BlowupConstruction.transSource_flat_off,TranscendentalNakamaye.Pending.BlowupConstruction.isOpen_transSource_flat_off,TranscendentalNakamaye.Pending.BlowupConstruction.differentiableAt_chart_pt,TranscendentalNakamaye.Pending.BlowupConstruction.differentiableOn_flat_off,TranscendentalNakamaye.Pending.BlowupConstruction.transSource_off_flat,TranscendentalNakamaye.Pending.BlowupConstruction.isOpen_transSource_off_flat,TranscendentalNakamaye.Pending.BlowupConstruction.differentiableOn_off_flat,TranscendentalNakamaye.Pending.BlowupConstruction.transSource_off_off,TranscendentalNakamaye.Pending.BlowupConstruction.isOpen_transSource_off_off,TranscendentalNakamaye.Pending.BlowupConstruction.differentiableOn_off_off,TranscendentalNakamaye.Pending.BlowupConstruction.isOpen_transSource_and_differentiableOn,TranscendentalNakamaye.Pending.BlowupConstruction.isClosed_centre,TranscendentalNakamaye.Pending.BlowupConstruction.exists_flattening,TranscendentalNakamaye.Pending.BlowupConstruction.flatK,TranscendentalNakamaye.Pending.BlowupConstruction.flatΦ,TranscendentalNakamaye.Pending.BlowupConstruction.flatΦ_spec,TranscendentalNakamaye.Pending.BlowupConstruction.mkFlat,TranscendentalNakamaye.Pending.BlowupConstruction.mkFlat_y₀,TranscendentalNakamaye.Pending.BlowupConstruction.mkFlat_Φ,TranscendentalNakamaye.Pending.BlowupConstruction.mkFlat_j,TranscendentalNakamaye.Pending.BlowupConstruction.exists_mem_flatSource,TranscendentalNakamaye.Pending.BlowupConstruction.exists_chart,TranscendentalNakamaye.Pending.BlowupConstruction.prefChart,TranscendentalNakamaye.Pending.BlowupConstruction.prefChart_mem,TranscendentalNakamaye.Pending.BlowupConstruction.mem_prefChart_source,TranscendentalNakamaye.Pending.BlowupConstruction.blowupCore,TranscendentalNakamaye.Pending.BlowupConstruction.topologicalSpace,TranscendentalNakamaye.Pending.BlowupConstruction.chartedSpace,TranscendentalNakamaye.Pending.BlowupConstruction.atlasHomeo,TranscendentalNakamaye.Pending.BlowupConstruction.atlasHomeo_toPartialEquiv,TranscendentalNakamaye.Pending.BlowupConstruction.isOpen_source_of_mem,TranscendentalNakamaye.Pending.BlowupConstruction.atlasHomeo_mem_atlas,TranscendentalNakamaye.Pending.BlowupConstruction.exists_eq_atlasHomeo,TranscendentalNakamaye.Pending.BlowupConstruction.isManifold,TranscendentalNakamaye.Pending.BlowupConstruction.FlatChart.isOpen_source,TranscendentalNakamaye.Pending.BlowupConstruction.FlatChart.homeo,TranscendentalNakamaye.Pending.BlowupConstruction.FlatChart.homeo_source,TranscendentalNakamaye.Pending.BlowupConstruction.FlatChart.homeo_apply,TranscendentalNakamaye.Pending.BlowupConstruction.continuousAt_proj_symm,TranscendentalNakamaye.Pending.BlowupConstruction.continuous_blowupProj,TranscendentalNakamaye.Pending.BlowupConstruction.exists_separating_of_common,TranscendentalNakamaye.Pending.BlowupConstruction.exists_separating_of_base_eq,TranscendentalNakamaye.Pending.BlowupConstruction.t2Space,TranscendentalNakamaye.Pending.BlowupConstruction.secondCountableTopology,TranscendentalNakamaye.Pending.BlowupConstruction.isComplexManifold,TranscendentalNakamaye.Pending.BlowupConstruction.contMDiffAt_proj_symm,TranscendentalNakamaye.Pending.BlowupConstruction.contMDiff_blowupProj}
\leanok
The maps \(\chi_c\) (\(c\) a flattening chart) and \(\chi_e\) (\(e\) a chart of \(X\)) form a holomorphic atlas of \(\tilde X\): every point of
\(\tilde X\) lies in the domain of one of them, and every coordinate change has an open domain and is holomorphic. For the topology it defines,
\(\tilde X\) is Hausdorff and second countable, hence a complex manifold of dimension \(n\), and \(\pi\) is holomorphic, with
\(\pi\circ\chi_c^{-1}=(\Phi\circ e_{y_0})^{-1}\circ\sigma_j\) and \(\pi\circ\chi_e^{-1}=e^{-1}\).
\end{theorem}
\begin{proof}
\leanok
\emph{Covering.} A point \((x,0)\), \(x\notin C\), lies in \(U_{e_x}\). For \(p=(y,L)\) with \(y\in C\), take a flattening at \(y\); since
\(\xi_p\notin T_yC\), Lemma~\ref{bex:tangent-flat}(2) gives a normal index \(j\) with \(d_y\tau_j(\xi_p)\ne0\), i.e.\ \(p\in U_c\).

\emph{Coordinate changes.} Between two charts \(\chi_e,\chi_{e'}\) it is a coordinate change of \(X\). From \(\chi_e\) to \(\chi_c\) it is
\(z\mapsto\chi_c(e^{-1}(z),0)\), whose coordinates are \(\tau_i(e^{-1}z)\) and \(\tau_i(e^{-1}z)/\tau_j(e^{-1}z)\), on the open set where
\(e^{-1}z\in D_\Phi\setminus C\) and \(\tau_j(e^{-1}z)\ne0\). From \(\chi_c\) to \(\chi_e\) it is \(w\mapsto e((\Phi\circ e_{y_0})^{-1}(\sigma_j w))\) on the
open set where \(w_j\ne0\) and \((\Phi\circ e_{y_0})^{-1}(\sigma_jw)\in e.\mathrm{source}\). From \(\chi_c\) to \(\chi_{c'}\), \(c'=(y_1,k',\Phi',j')\) with
flat coordinates \(\tau'\): on \(U_c\), \(\tau_i'^\flat/\tau_{j'}'^\flat=(\tau_i'^\flat/\tau_j^\flat)/(\tau_{j'}'^\flat/\tau_j^\flat)\), so in the coordinates \(w\) of
\(\chi_c\) the coordinate change is
\[
 w\longmapsto\Bigl(\tau_i'\bigl((\Phi\circ e_{y_0})^{-1}(\sigma_jw)\bigr)\ (i<k'\text{ or }i=j'),\quad G_{\tau_i'}(w)/G_{\tau_{j'}'}(w)\ (k'\le i\ne j')\Bigr),
\]
on the set where \((\Phi\circ e_{y_0})^{-1}(\sigma_jw)\in D_{\Phi'}\) and \(G_{\tau_{j'}'}(w)\ne0\). The functions \(\tau_i'\), \(i\ge k'\), are holomorphic on
\(D_{\Phi'}\) and vanish on \(C\cap D_{\Phi'}\), so by Lemma~\ref{bex:ratio-holomorphic} the \(G\) are holomorphic, the domain is open and the change
is holomorphic.

So the charts generate a topology in which each \(U\) is open and each \(\chi\) a homeomorphism onto an open subset of \(\mathbf C^n\), and the atlas
is holomorphic. The formulas for \(\pi\circ\chi^{-1}\) (Lemma~\ref{bex:monomial-chart}) show that \(\pi\) is continuous and holomorphic.

\emph{Hausdorff.} Let \(p\ne q\). If \(\pi(p)\ne\pi(q)\), pull back disjoint neighbourhoods in \(X\) by \(\pi\). Otherwise \(\pi(p)=\pi(q)=y\in C\)
(over a point off \(C\) there is one point). Fix a flattening at \(y\) and normal indices \(j,l\) with \(p\in U_{(\Phi,j)}\), \(q\in U_{(\Phi,l)}\). If
\(q\in U_{(\Phi,j)}\) or \(p\in U_{(\Phi,l)}\), the two points lie in one chart domain and are separated by the chart. Otherwise
\(\tau_j^\flat(q)=0=\tau_l^\flat(p)\), so \(l\ne j\), \(|\chi_{(\Phi,j)}(p)_l|=0\) and \(|\chi_{(\Phi,l)}(q)_j|=0\); the open sets
\(\{r\in U_{(\Phi,j)}:|\chi_{(\Phi,j)}(r)_l|<1\}\ni p\) and \(\{r\in U_{(\Phi,l)}:|\chi_{(\Phi,l)}(r)_j|<1\}\ni q\) are disjoint, since on
\(U_{(\Phi,j)}\cap U_{(\Phi,l)}\) the product \(\chi_{(\Phi,j)}(r)_l\,\chi_{(\Phi,l)}(r)_j=(\tau_l^\flat/\tau_j^\flat)(\tau_j^\flat/\tau_l^\flat)\) is \(1\).

\emph{Second countability.} Countably many charts of \(X\) cover \(X\), and since \(C\) is closed in the second countable space \(X\),
countably many flat domains \(D_{\Phi_m}\) cover \(C\). The corresponding countably many chart domains cover \(\tilde X\) (as in the covering step), and
each is homeomorphic to an open subset of \(\mathbf C^n\).
\end{proof}

\begin{lemma}[The exceptional divisor]
\label{bex:proj-invertible}
\uses{bex:atlas,bex:hadamard}
\lean{TranscendentalNakamaye.Pending.BlowupConstruction.exists_coord_ne_zero_of_mem_nhds,TranscendentalNakamaye.Pending.BlowupConstruction.holomorphicAt_comp_blowupProj,TranscendentalNakamaye.Pending.BlowupConstruction.FlatChart.holomorphicAt_coord,TranscendentalNakamaye.Pending.BlowupConstruction.FlatChart.τ_base_eq,TranscendentalNakamaye.Pending.BlowupConstruction.isInvertible_pullback_blowupProj}
\leanok
The ideal \(\pi^*\mathcal I_C\) is invertible: on \(U_c\) it is generated by \(\tau_j\circ\pi\), which is the coordinate \(w_j\) of \(\chi_c\), and on
\(\pi^{-1}(X\setminus C)\) it is the unit ideal.
\end{lemma}
\begin{proof}
\leanok
Off \(C\) the constant \(1\) lies in \(\mathcal I_C\) (it vanishes on \(C\) near points off the closed set \(C\)). On \(U_c\),
\(\tau_j\in\mathcal I_C\) on \(D_\Phi\), so \(\tau_j\circ\pi\) lies in the pullback. Conversely every \(b\in\mathcal I_{C,\pi(p)}\) is
\(\sum_{i\ge k}\beta_i\tau_i\) near \(\pi(p)\) (Lemma~\ref{bex:hadamard}), and \(\tau_i\circ\pi\circ\chi_c^{-1}=\sigma_j(w)_i\) is \(w_iw_j\) or \(w_j\) for
\(i\ge k\); so \(b\circ\pi\) is a holomorphic multiple of \(w_j\). Finally \(w_j\) vanishes on no nonempty open subset of the chart.
\end{proof}

\begin{lemma}[Uniqueness of lifts]
\label{bex:lift-unique}
\uses{bex:atlas}
\lean{TranscendentalNakamaye.Pending.BlowupConstruction.exists_not_mem_of_isOpen,TranscendentalNakamaye.Pending.BlowupConstruction.lift_eqOn}
\leanok
Let \(g\colon W\to X\) be a holomorphic map from a complex manifold with \(g^*\mathcal I_C\) invertible. Then \(g^{-1}(X\setminus C)\) is dense
in \(W\), and two continuous maps \(h_1,h_2\colon W'\to\tilde X\) on an open set \(W'\subset W\) with \(\pi\circ h_1=\pi\circ h_2=g\) on \(W'\)
coincide.
\end{lemma}
\begin{proof}
\leanok
Suppose a nonempty open \(W_1\) satisfies \(g(W_1)\subset C\). Let \(w\in W_1\) and \(u\) a generator of \(g^*\mathcal I_C\) near \(w\). Then \(u\) lies
in \(g^*\mathcal I_C\) at \(w\): near \(w\), \(u=\sum_lc_l\,(b_l\circ g)\) with \(b_l\in\mathcal I_{C,g(w)}\). Each \(b_l\) vanishes on \(C\) near
\(g(w)\), so \(u\) vanishes on a nonempty open subset of \(W_1\) near \(w\), which the generator does not. For the second statement, a point of
\(\tilde X\) over \(x\notin C\) is \((x,0)\), so \(h_1=h_2\) on \(W'\cap g^{-1}(X\setminus C)\), a dense subset of \(W'\); and
\(\{h_1\ne h_2\}\) is open since \(\tilde X\) is Hausdorff (Theorem~\ref{bex:atlas}). An open set missing a dense set is empty.
\end{proof}

\begin{lemma}[Local existence of lifts]
\label{bex:lift-local}
\uses{bex:atlas,bex:hadamard}
\lean{TranscendentalNakamaye.Pending.BlowupConstruction.contMDiffAt_of_mem_chart,TranscendentalNakamaye.Pending.BlowupConstruction.exists_local_lift_of_not_mem,TranscendentalNakamaye.Pending.BlowupConstruction.exists_local_lift_of_mem,TranscendentalNakamaye.Pending.BlowupConstruction.exists_local_lift}
\leanok
Let \(g\colon W\to X\) be a holomorphic map from a complex manifold with \(g^*\mathcal I_C\) invertible, and \(w_0\in W\). There are an open
\(W_0\ni w_0\) and a holomorphic \(h\colon W_0\to\tilde X\) with \(\pi\circ h=g\) on \(W_0\).
\end{lemma}
\begin{proof}
\leanok
If \(g(w_0)\notin C\), take \(h(w)=(g(w),0)=\chi_e^{-1}(e(g(w)))\) on \(g^{-1}(e.\mathrm{source}\setminus C)\). Let \(g(w_0)\in C\); take a flattening
at \(g(w_0)\) with flat coordinates \(\tau_i\), and a generator \(u\) of \(g^*\mathcal I_C\) near \(w_0\). The functions \(f_i=\tau_i\circ g\), \(i\ge k\),
lie in \(g^*\mathcal I_C\), so \(f_i=a_iu\) near \(w_0\) with \(a_i\) holomorphic. And \(u\) lies in \(g^*\mathcal I_C\):
\(u=\sum_lc_l\,(b_l\circ g)\) near \(w_0\) with \(b_l\in\mathcal I_{C,g(w_0)}\), and \(b_l=\sum_{i\ge k}\beta_{li}\tau_i\) near \(g(w_0)\)
(Lemma~\ref{bex:hadamard}), so \(u=\sum_{i\ge k}\gamma_if_i=u\sum_{i\ge k}\gamma_ia_i\) near \(w_0\) with \(\gamma_i\) holomorphic. If all
\(a_i(w_0)\) vanished, \(1-\sum\gamma_ia_i\) would be nonzero near \(w_0\) and \(u\) would vanish on a neighbourhood of \(w_0\), which a generator
does not. So \(a_j(w_0)\ne0\) for some \(j\ge k\). Near \(w_0\) let \(z(w)_i=\tau_i(g(w))\) for \(i<k\) or \(i=j\), and
\(z(w)_i=a_i(w)/a_j(w)\) for \(k\le i\ne j\). Since \(\tau_i(g(w))=a_iu=(a_i/a_j)\,\tau_j(g(w))\), \(\sigma_j(z(w))=\Phi(e_{y_0}(g(w)))\in\Phi.\mathrm{target}\);
so \(h=\chi_{(\Phi,j)}^{-1}\circ z\) is defined and holomorphic near \(w_0\), and \(\pi\circ h=(\Phi\circ e_{y_0})^{-1}\circ\sigma_j\circ z=g\)
(Theorem~\ref{bex:atlas}).
\end{proof}

\begin{theorem}[The blowup exists]
\label{bex:blowup-exists}
\uses{bex:atlas,bex:proj-invertible,bex:lift-unique,bex:lift-local,ana:blowups}
\lean{TranscendentalNakamaye.Pending.BlowupConstruction.exists_lift,TranscendentalNakamaye.Pending.BlowupConstruction.isBlowupAlong_blowupProj}
\leanok
Let \(C\) be a nowhere dense closed submanifold of a complex manifold \(X\). Then \(\pi\colon\tilde X\to X\) is a blowup of \(X\) along \(C\) in
the sense of Definition~\ref{ana:blowups}.
\end{theorem}
\begin{proof}
\leanok
\(\tilde X\) is a complex manifold of dimension \(n\) and \(\pi\) is holomorphic (Theorem~\ref{bex:atlas}), and \(\pi^*\mathcal I_C\) is invertible
(Lemma~\ref{bex:proj-invertible}). Let \(g\colon W\to X\) be holomorphic, \(W\) a complex manifold of any dimension, with \(g^*\mathcal I_C\)
invertible. Every point of \(W\) has a neighbourhood with a holomorphic lift of \(g\) (Lemma~\ref{bex:lift-local}); two such lifts agree on the
overlap of their domains (Lemma~\ref{bex:lift-unique}), so they glue to a holomorphic \(h\colon W\to\tilde X\) with \(\pi\circ h=g\), and \(h\) is
unique by Lemma~\ref{bex:lift-unique}.
\end{proof}
